%==========================================================
% Synapse AI (TrustMed AI): A Neuro-Symbolic, Multimodal
% Clinical Decision Support System with Adversarial Safety
%
% main3.tex
% Humanized voice, real evaluation tables, figure placeholders.
% Compile with pdfLaTeX (Overleaf: asmeconf.cls is included).
%==========================================================

\documentclass[twocolumn,10pt,main=USenglish]{asmeconf}

\usepackage{graphicx}
\usepackage{amsmath}
\usepackage{booktabs}
\usepackage{multirow}
\usepackage{siunitx}
\sisetup{group-separator={,}, group-minimum-digits=4}
\usepackage{hyperref}
\usepackage{url}
\usepackage{xcolor}
\usepackage{array}
\usepackage{caption}
\usepackage{float}
\usepackage{listings}

\newcolumntype{L}[1]{>{\raggedright\arraybackslash}p{#1}}
\captionsetup{justification=raggedright,singlelinecheck=false}
\raggedbottom
\emergencystretch=2em

\hypersetup{
    colorlinks=true,
    linkcolor=black,
    citecolor=black,
    urlcolor=blue
}

%-----------------------------------------------------------
% Paper metadata (ASME fields)
%-----------------------------------------------------------
\ConfName{FSE 570 Capstone Project Report}
\ConfAcronym{FSE570}
\ConfDate{April 30, 2026}
\ConfCity{Tempe, AZ, USA}
\PaperNo{SYNAPSE-2026-001}
\title{Synapse AI: A Neuro-Symbolic, Multimodal Clinical\\
Decision Support Pipeline with an Adversarial\\
Safety Critic}

\SetAuthors{%
  Chandana Pulikanti\affil{1},\;
  Mansi Nandkar\affil{1},\;
  Pragyan Jyoti Borthakur\affil{1},\;
  Pranav Sunil\affil{1},\;
  Pritish Kannan\affil{1}%
}
\SetAffiliation{1}{School of Computing and Augmented Intelligence,
Arizona State University, Tempe, AZ, USA.\\
Emails: 
{cpulikan@asu.edu}, \texttt{mnandkar@asu.edu},
\texttt{pborthak@asu.edu}, \texttt{psunil@asu.edu}, \texttt{pkannan7@asu.edu}}

\keywords{Clinical Decision Support, Neuro-Symbolic AI, Knowledge
Graphs, Retrieval-Augmented Generation, Medical Imaging,
LLM-as-Judge, Adversarial Safety}

\begin{document}
\maketitle

%===========================================================
\begin{abstract}
Large language models deployed in clinical settings exhibit
confident hallucination of dosing, contraindications, and imaging
findings. This paper presents Synapse AI, a clinical decision support pipeline that addresses these
failure modes through architectural decomposition:  deterministic graph and vector retrieval, multimodal vision through
MedGemma~\cite{medgemma2024}, and
an adversarial second-LLM safety critic with authority to override
the draft. We evaluate the pipeline across five suites. Clinical-text suites
(Retrieval QA, A/B Latency, Adversarial Stress, Golden Dataset)
draw from the MIMIC-IV Demo cohort, including a polypharmacy chart
with up to \num{88} active medications. The Blind Vision suite
draws from held-out MIMIC-CXR studies. The ROCO/PMC subset is
retained as development data and is not used to compute any
reported metric. The pipeline achieves a macro-average entity-level precision of
${\sim}\num{92}$~percent across three temperature-0 runs on an
expanded Golden Dataset (\num{25} queries across five MIMIC-IV Demo
patients) with the safety critic layer catching severe medical hallucinations under the
Section~\ref{sec:metric-defs} definition. In the adversarial
suite, seven of eight cases scored \num{9}/\num{10} groundedness,
and the eighth was blocked through a critic-triggered override.
The pipeline incurs a \num{9.04}-second mean safety overhead, an
overhead we argue is acceptable in safety-critical clinical
decision support.
\end{abstract}

%===========================================================
\section{Introduction}
Clinical decision support requires assistants that ground every
response in the patient record, refuse confidently when evidence is
absent, and exhibit predictable failure modes~\cite{topol2019deep,
rajpurkar2022ai}. Single-pass LLMs satisfy the first two requirements
inconsistently and the third rarely~\cite{ji2023survey}. A model that summarizes a chart correctly in one session may overlook an active clopidogrel prescription before recommending warfarin in the next, because a single forward pass is required to recall facts,
reason over them, generate prose, and self-verify simultaneously.
Architectural decomposition is therefore necessary to separate
these responsibilities~\cite{moor2023foundation}.

Synapse AI implements this decomposition. The pipeline routes each
query through three independent retrievers (patient record,
knowledge graph, vector store), fuses their output into a grounded
context, generates a draft response with an instruction-tuned model,
and submits the draft to a separate critic model for adversarial
review. If the critic identifies an inconsistency, the draft is
replaced by a flagged override before the response reaches the user.

The contributions of this paper are as follows.
(1) A neuro-symbolic orchestrator (\texttt{trustmed\_brain}) that
runs parallel retrieval across three context channels (patient
record, knowledge graph, vector store), assembled by a Context
Binding step, and exposes the result through a streaming FastAPI
backend.
(2) A multimodal vision pipeline that pairs MedGemma analysis with
BiomedCLIP cross-reference validation against a labeled MIMIC-CXR
anchor bank to flag findings unsupported by similar historical
cases.
(3) An evaluation across five suites (retrieval QA, A/B latency
with and without the critic, adversarial stress, blind vision, and
Golden Dataset entity-level validation) reported on numbers drawn
from the production stack.

\subsection{Related Work}
\label{sec:related}

Hallucinated diagnoses,
fabricated doses, and ungrounded imaging are the dominant failure
mode of LLMs in clinical workflows~\cite{ji2023survey,
pal2023medhalt}. MedPaLM~\cite{singhal2023large} reaches strong
USMLE accuracy but still fabricates on open-ended chart
queries~\cite{omiye2023racial}.

Neuro-symbolic systems pair LLMs with symbolic constraints so that
outputs trace to verifiable sources~\cite{sheth2023neurosymbolic}.
Medical pipelines layer this with a structured knowledge graph
reasoning over clinical ontologies~\cite{rotmensch2017graph}.

Dense vector retrieval with cross-encoder rerank is the open-domain
default~\cite{lewis2020rag, reimers2019sbert}. Graph-RAG
variants~\cite{edge2024local} trade dense coverage for explicit
multi-hop grounding via typed edges.

LLM-as-judge auditing
spans evaluation~\cite{zheng2023judge} and runtime safety;
Adversarial second-LLM auditors have been proposed for
red-teaming~\cite{perez2022red} and clinical post-hoc
review~\cite{karabacak2023mirroring}.

%===========================================================
\section{Data Sources and Data Analysis}
\label{sec:data}

Synapse AI integrates two heterogeneous datasets and three curated
knowledge sources. Table~\ref{tab:data} summarizes scale and role.

\begin{table}[t]
\centering
\caption{Datasets and knowledge sources.}
\label{tab:data}
\small
\begin{tabular}{@{}L{0.24\linewidth}L{0.30\linewidth}L{0.36\linewidth}@{}}
\toprule
\textbf{Source} & \textbf{Scale} & \textbf{Role} \\
\midrule
Kaggle Disease--Symptom--\newline Medicine corpus
 & ${\sim}\num{1676}$ diseases, \num{3920} symptoms, \num{3867} medicines
 & ChromaDB collections + Neo4j baseline graph. \\
\addlinespace
ROCO / PMC medical image subset~\cite{pelka2018roco}
 & \num{1800} train, \num{200} test
 & Development image corpus; not used as anchor bank. \\
\addlinespace
UMLS Metathesaurus~\cite{bodenreider2004umls}
 & ${\sim}\num{4.6}$M concepts
 & Concept normalization and graph enrichment. \\
\addlinespace
MIMIC-IV Demo~\cite{johnson2023mimiciv}
 & \num{12} patients, \num{600}+ records
 & Clinical-text suites (de-identified, public). \\
\addlinespace
MIMIC-CXR (anchor bank)~\cite{johnson2019mimiccxr}
 & ${\sim}\num{50000}$ images
 & Visual-RAG bank, BiomedCLIP-embedded~\cite{zhang2024biomedclip}; test cases excluded (Section~\ref{sec:setup}). \\
\bottomrule
\end{tabular}
\end{table}

\subsection{Cleaning, Preprocessing, and Integration}
Ingestion runs four steps. Schema validation against a declarative
config rejects malformed rows up front. Entity normalization maps
strings to UMLS CUIs. Deduplication combines exact match with
embedding-similarity clustering. Each surviving row dual-writes to
ChromaDB (embeddings) and Neo4j (symbolic relations). Medical
images go through a compound-figure detector first; multi-panel
layouts get split before embedding so panel provenance survives
into retrieval. The graph that comes out has explicit disease,
symptom, and drug nodes; the orchestrator can query it without
asking the LLM to recall medical facts.

\subsection{Provenance, Privacy, and Threat Model}
\label{sec:provenance}
All patient identifiers (\texttt{10040025}, \texttt{10002428},
\texttt{10014354}, \texttt{10015860}, \texttt{10016742},
\texttt{10023239}, \texttt{10035631}) resolve against the MIMIC-IV
Demo cohort: publicly licensed, fully de-identified, no
PHI~\cite{johnson2023mimiciv}. Clinical-text suites (Retrieval QA,
A/B Latency, Adversarial Stress, Golden Dataset) run against
MIMIC-IV Demo; Blind Vision runs against held-out MIMIC-CXR
studies~\cite{johnson2019mimiccxr}. A developer-only synthetic
patient store exists as an offline sandbox and backs no reported
result.

Real-PHI ingestion is opt-in behind a credentialed path. Uploads
go through a canonical root with symlink and path-traversal
checks. Outbound traffic uses three isolated surfaces: OpenRouter
(draft + critic), Vertex AI MedGemma (vision), UMLS REST. An
operator can redact, swap, or disable any one without touching the
orchestrator. Institutional deployment reuses that boundary inside
a VPC-bound reverse proxy with provider-side audit logging,
keeping the fail-closed posture intact.

%===========================================================
\section{Methodology and Application}
\label{sec:method}

The pipeline is organized as a multi-layer neuro-symbolic stack
(Table~\ref{tab:layers}). Layers~4 (Drug Safety Engine) and~7
(Visual-RAG Consistency Check) are fully deterministic. Layers~6
(Model) and~8 (Safety Critic) are model-based. The Pipeline
Router (L1), parallel retrieval (L2/L3), and Context Binding (L5)
sit between the two: deterministic in their dispatch logic,
data-driven in what they return. Deterministic layers constrain
what the model layers see and emit; the Safety Critic provides
post-hoc audit, not a deterministic guarantee.

\begin{table}[H]
\centering
\caption{Multi-layer architecture overview.}
\label{tab:layers}
\small
\begin{tabular}{@{}cL{0.20\linewidth}L{0.56\linewidth}@{}}
\toprule
\textbf{L} & \textbf{Layer} & \textbf{Role}\\
\midrule
1 & Pipeline Router & Routes each query to the right retrieval modules (text, vision, PDF) by input type.\\
2 & Fact + Graph & SQLite patient context ($\phi_P$) is fetched first because the prompt depends on it; Neo4j Cypher ($\phi_G$) and Vector-RAG ($\phi_V$, Layer~3) are then issued concurrently via \texttt{asyncio.gather}.\\
3 & Vector-RAG + Rerank & ChromaDB dense retrieval ($\phi_V$, $N{=}32$) + cross-encoder rerank ($k{=}8$); launched in parallel with $\phi_G$.\\
4 & Drug Safety Engine & Deterministic set-intersection over meds; injects \texttt{[HIGH RISK ALERT]} pre-synthesis. Bypasses the LLM.\\
5 & Context Binding & Assembles $\phi_P \oplus \phi_G \oplus \mathrm{Rerank}_k(\mathcal{D}(q))$ plus vision and chat history into one prompt.\\
6 & Model Layer & Nemotron (text) or MedGemma (vision) drafts the response.\\
7 & Visual-RAG Check & Post-synthesis cosine match against the BiomedCLIP/MIMIC-CXR anchor; overrides if vision contradicts anchor.\\
8 & Safety Critic & Gemma~3 audits the draft against \texttt{[PatientContext]}. Safe: pass. Unsafe + correction: override. Ambiguous / timeout / error: keep draft.\\
\bottomrule
\end{tabular}
\end{table}

\subsection{System Architecture}
Figure~\ref{fig:arch} shows the end-to-end pipeline. A Next.js 15
frontend sends text and vision inputs to a FastAPI backend. A
\textbf{Pipeline Router} classifies each incoming query and activates
the appropriate retrieval modules. Three context channels run in
parallel: patient fact retrieval from SQLite ($\phi_P$), graph
retrieval from Neo4j ($\phi_G$), and vector retrieval from ChromaDB
with cross-encoder rerank ($\phi_V$). The \textbf{Drug Safety Engine}
runs a deterministic set-intersection over the active medication list
and injects hard-lock alerts into the context before synthesis.
\textbf{Context Binding} assembles all channels into a single grounded
prompt (Eq.~\eqref{eq:fuse}). The \textbf{Model Layer} (Nemotron for
text, MedGemma for vision) generates a draft response. A
\textbf{Visual-RAG Consistency Check} then compares draft claims
against the BiomedCLIP anchor bank deterministically. Finally, an
adversarial \textbf{Safety Critic} (Gemma~3, isolated from the
synthesizer) audits the draft and may override it before the response
is streamed to the user.

\begin{figure*}
  \centering
  \includegraphics[width=\linewidth]{figures/FINAL - Synapse AI.png}
  \caption{Synapse AI end-to-end architecture.}
  \label{fig:arch}
\end{figure*}

\subsection{Neuro-Symbolic Orchestration}
The orchestrator (\texttt{trustmed\_brain}) fuses three
independent context functions before generation. Let $q$ be the user
query and $p$ the active patient identifier. Each context function
returns an ordered list of passages. The fused context is the
ordered concatenation
\begin{equation}
\mathcal{C}(q,p) \;=\; \phi_{P}(p) \;\oplus\; \phi_{G}(q,p) \;\oplus\;
                       \mathrm{Rerank}_k\!\left(\mathcal{D}(q)\right),
\label{eq:fuse}
\end{equation}
where $\oplus$ denotes order-preserving concatenation and the three
operands are $\phi_P(p)$, the patient context drawn from
\texttt{mimic\_demo.db}; $\phi_G(q,p)$, a Cypher query against the
Neo4j graph that takes $p$ only as a read-only cohort filter (no
shared mutable state with $\phi_V$);
and $\phi_V(q)$, a dense top-$N$ pull from ChromaDB. The reranker is
a cross-encoder $s_\psi(q,d)$ that scores query and candidate
jointly,
\begin{equation}
\begin{aligned}
\mathcal{D}(q) &:= \phi_V(q), \quad |\mathcal{D}(q)| = N, \\
\mathrm{Rerank}_k\!\left(\mathcal{D}(q)\right) &= \bigl[\, d_{(1)},\ldots,d_{(k)} \,\bigr], \\
s_\psi(q, d_{(1)}) &\geq \cdots \geq s_\psi(q, d_{(k)}).
\end{aligned}
\label{eq:rerank}
\end{equation}
which returns an ordered list of length $k$ (ties broken by the
upstream dense-retrieval rank). In production $k{=}8$, $|\mathcal{D}(q)|
{=}\num{32}$.

Per-turn loop. Resolve the patient, pull $\phi_P$ first because
the prompt depends on it. Then $\phi_G$ and $\phi_V$ go out
concurrently under one \texttt{asyncio.gather}; $\phi_G$ uses the
patient identifier as a read-only filter, so the two calls share
no mutable state. Eq.~\eqref{eq:fuse} fuses the three contexts;
length-aware packing trims to the prompt budget. The draft $r_0$
streams token-by-token over SSE for partial-answer rendering.
After the stream closes, the critic $\sigma$ reads $r_0$ against
$\phi_P$ and returns one of three outcomes: emit $r_0$, emit a
critic-revised $r$, or request a regeneration with an evidence
nudge. Running the audit after the stream keeps time-to-first-
token free of safety overhead while still gating the final
response before it leaves the backend.

\subsection{Adversarial Safety Critic}
Two agents, distinct roles. Synthesizer:
\texttt{nvidia/nemotron-3-nano-30b-a3b:free} at $T{=}\num{0.1}$,
drafts the response. Auditor:
\texttt{google/gemma-3n-e4b-it:free} at $T{=}\num{0.0}$, served
through OpenRouter, blind to the user query. The auditor sees only
the draft and the raw \texttt{[PatientContext]} arrays. Three
checks run in order. \emph{Hallucination}: any ICD-10 term in the
draft missing from the Layer~2 SQL payload. \emph{Omission}: any
Drug Safety Engine alert (Layer~4) the draft skipped.
\emph{Adversarial payload}: jailbreak patterns. A failed check
drops the draft and substitutes a flagged override. Timeout above
\num{30}~s, ambiguous verdict, or any exception keeps the draft
intact and logs the failure; nothing is dropped silently. Gate, not
silent rewriter. Every override is logged.

\subsection{Drug Safety Engine (Layer 4)}
Drug safety skips the LLM. Set-intersection over the active
medication list against curated pharmacological knowledge bases.
A hit on warfarin + ibuprofen, CYP1A2 inhibitor + narrow-therapeutic-
index agent, or QT-prolonging combination injects
\texttt{[HIGH RISK ALERT]} into the prompt context. Synthesizer is
told to surface the alert verbatim; Layer~8 overrides if missing.
Two redundant lines of defense, no token-level decode.

\subsection{Multimodal Vision Pipeline}
Each uploaded image runs through four phases. Phase~1: MedGemma
returns a structured report tagging every observation as
\texttt{[HIGH]} or \texttt{[LOW]} confidence. Phase~2: build a
conservative search query from modality, region, and only the
\texttt{[HIGH]}-confidence findings. \texttt{[LOW]} lines never
seed retrieval; that keeps uncertain wording from polluting the
downstream prompt. Phase~3: run the query against the
\texttt{diseases} ChromaDB collection with cross-encoder rerank
(top-$k{=}3$, normalized score $\geq 0.30$). Only guidelines past
the relevance bar reach the prompt. Phase~4: BiomedCLIP visual
search against the MIMIC-CXR anchor bank, keeping matches above
\num{70}~percent cosine similarity. A SHA-256 hash of the image
bytes acts as the cache key, so a re-uploaded scan or a follow-up
question returns the cached pipeline result without re-invoking
MedGemma. Compound multi-panel figures are processed as one unit
today; splitting is on the Section~\ref{sec:future} list.

\subsection{Visual-RAG Anchoring and Cross-Reference Validation}
A deterministic cross-reference step sits between vision output
and the Safety Critic. BiomedCLIP~\cite{zhang2024biomedclip}
encodes the query image and pulls the top \num{10} nearest
neighbors from the ${\sim}\num{50000}$-image MIMIC-CXR anchor
bank~\cite{johnson2019mimiccxr} at cosine similarity
$\geq \num{0.60}$. The agent maps each \texttt{[HIGH]}-confidence
finding to canonical CheXpert labels through a hand-curated alias
table; ``opacification'' $\to$ \texttt{lung opacity},
``collapsed lung'' $\to$ \texttt{pneumothorax}. Verdict logic.
A finding is \emph{corroborated} when the matched condition
appears in $\geq \num{30}$~percent of neighbors,
\emph{uncorroborated} otherwise. Uncorroborated findings ride into
the report flagged as possible hallucinations; the Safety Critic
treats the flag as extra context. The agent also surfaces
omissions: conditions present in $\geq \num{50}$~percent of
similar images that the vision model never mentioned. Below three
qualifying MIMIC-CXR neighbors, the step reports insufficient data
rather than emitting a verdict.
\emph{Leakage control.} The three blind-vision test cases
(\texttt{test\_1963}, \texttt{test\_1852}, \texttt{test\_1806})
and their near-duplicates were dropped from the anchor bank before
evaluation; see Section~\ref{sec:vision-results}.

%===========================================================
\section{Experimental Setup}
\label{sec:setup}

We ran five evaluation suites against the production pipeline.

\subsection{Operational Metric Definitions}
\label{sec:metric-defs}
\emph{Severe medical hallucination} = a clinically actionable claim
(diagnosis, drug, dose) in the response that is not in the
patient's \texttt{true\_diagnoses} or \texttt{true\_medications}.
Bold markdown headers and vital discrepancies do not count; both
get filtered before scoring. \emph{F1} = harmonic mean of precision
and recall on medication and diagnosis entities,
$F_1 = 2\cdot\tfrac{P\cdot R}{P+R}$. Reported per query and as a
macro-average. \emph{Safety recall} = share of a patient's
\texttt{known\_contraindications} surfaced in the response by
keyword coverage at $\geq \num{50}$~percent; N/A where no
contraindications are defined. \emph{Vital discrepancy} = a
numeric vital outside the per-vital $\pm$tolerance window. Tracked
separately, excluded from the precision denominator.

\textbf{Retrieval QA:} \num{34} clinician-style queries across seven
chronic conditions: asthma, cholesterol, diabetes, heart disease,
hypertension, migraine, stroke. Each response gets three scores;
keyword overlap against a clinician-written reference, cosine
similarity against the same reference, and end-to-end latency.
Mix: \num{15} single-turn, \num{3} edge cases, \num{16}
conversational follow-ups for context rewrite.

\textbf{A/B latency:} Five complex multi-drug queries through two
configurations. Pipeline~A monkey-patches
\texttt{use\_llm\_safety}=\texttt{False} and returns the
single-pass draft. Pipeline~B is the production stack with the
critic on. Mean and P95 latency reported.

\textbf{Adversarial stress:} Eight escalating cases. Standard
chart consult through a panoramic mandible X-ray queried for
pneumonia, a hip MRI on a steroid-treated patient, a prompt
injection (``ignore all previous instructions''), an off-topic Star
Wars summary. Gemma~3 judge scores each output for groundedness
against the patient record.

\textbf{Blind vision:} Unannotated MIMIC-CXR studies were fed
to the vision pipeline with a generic prompt: ``Interpret this scan
blindly and detail any pathological findings.'' Outputs were graded
by the same Gemma~3 judge against MIMIC-CXR ground-truth labels for
the corresponding studies. Cases \texttt{test\_1963},
\texttt{test\_1852}, and \texttt{test\_1806} are MIMIC-CXR studies
held out from the BiomedCLIP anchor bank.

\textbf{Golden Dataset validation:} A separate suite replaces
LLM-as-judge with rigid entity-level matching. Verified clinical
entities (\texttt{true\_diagnoses}, \texttt{true\_medications},
\texttt{true\_vitals}, \texttt{known\_contraindications}) are
extracted from the MIMIC-IV source tables for each target patient
and serialized into a versioned golden dataset. The validator runs
two extraction phases against every response. Phase~1 scans the
full text against the patient's known medication and diagnosis
vocabulary using fuzzy boolean matching (Levenshtein similarity
$\geq 0.80$). Phase~2 scores bold markdown spans
(\texttt{**term**}) against the same vocabulary, filtered through
the \texttt{CLINICAL\_SECTION\_HEADERS} vocabulary so that
structural headings such as \textbf{Glycemic Optimization} or
\textbf{Underlying Conditions} are not mistaken for drug names.
Vital signs are extracted via regex and checked against per-vital
tolerance windows; out-of-window values are classified as vital
discrepancies (Section~\ref{sec:metric-defs}) and excluded from the
precision denominator. The suite runs \num{25} queries across five
patients (\texttt{10040025}, \texttt{10002428}, \texttt{10014354},
\texttt{10015860}, \texttt{10023239}); each patient is tested on
five query types: general assessment, drug-safety add-on,
condition-specific management, discharge planning, and single-drug
safety check. Reported metrics are precision, recall, F1, and
safety recall as defined in Section~\ref{sec:metric-defs}.

\subsection{LLM Judge Rubric}
\label{sec:judge-rubric}
Gemma~3 scores adversarial-stress and blind-vision outputs on four
0--10 axes, equally weighted. The judge never sees the user query;
inputs are the response plus patient context or ground-truth
label. \emph{Faithfulness}: $-2$ per fabricated entity.
\emph{Completeness}: $-2$ per material omission. \emph{Safety}:
collapses to \num{0} on any unflagged unsafe recommendation.
\emph{Calibration}: hedging on uncertain claims. Axis scores get
averaged, rounded, and reported as the headline x/\num{10} in
Section~\ref{sec:adv-results} and Table~\ref{tab:vision}.
Overrides are categorical, not numeric
(Section~\ref{sec:metric-defs}).

%===========================================================
\section{Results}
\label{sec:results}

\subsection{Retrieval QA across Conditions}
Figure~\ref{fig:retrieval} reports per-topic retrieval performance
across the \num{34}-query suite. Heart disease and migraine score
highest on cosine similarity (\num{0.632} and \num{0.642}). Stroke
falls to \num{0.390} despite hitting every keyword, traced to
reference answers favoring outcome-oriented phrasing while the
pipeline answers in mechanism-oriented phrasing. Asthma is the only
topic with reduced keyword coverage (\num{80}~percent), driven by
two queries that reference a brand-name inhaler absent from the
corpus.

Table~\ref{tab:diff} stratifies the same suite by difficulty. The
difficulty label is taken from the test fixture and reflects how
many hops of reasoning the reference answer assumes. Conversational
context rewrite (the rate at which a follow-up like ``and what about
his beta blocker?'' is correctly rewritten with the prior patient
context) is \num{16} of \num{16}.

\begin{table}[H]
\centering
\caption{Difficulty stratification and aggregate retrieval metrics.}
\label{tab:diff}
\small
\begin{tabular}{@{}lcc@{}}
\toprule
\textbf{Tier} & \textbf{Keyword Match} & \textbf{Avg. Sim.}\\
\midrule
Easy   & \num{96}\% & \num{0.592} \\
Medium & \num{95}\% & \num{0.443} \\
Hard   & \num{90}\% & \num{0.436} \\
\midrule
Conversational rewrite & \multicolumn{2}{c}{\num{16}/\num{16}} \\
Mean similarity (all)  & \multicolumn{2}{c}{\num{0.500} $\pm$ \num{0.138}}\\
Mean keyword match (all) & \multicolumn{2}{c}{\num{93}\%}\\
Mean latency (all)     & \multicolumn{2}{c}{\num{28423} $\pm$ \num{5231}~ms}\\
\bottomrule
\end{tabular}
\end{table}


\textbf{Cosine divergence is a register mismatch.} Stroke at
\num{0.390}, Cholesterol at \num{0.402}, both at \num{100}\%
keyword match (Figure~\ref{fig:retrieval}). The reference answers
phrase recommendations in outcome terms (``lower bad
cholesterol''); the graph returns mechanism terms (``HMG-CoA
reductase inhibition''). Same clinical content, different
register. Cosine penalizes the vocabulary swap; keyword match,
treated as set membership, sees the entities are present.

The graph and vector layers retrieved the right clinical entities
and the synthesizer surfaced them. The metric gap comes from
cosine, not from retrieval failure.

\begin{figure}[H]
  \centering
  \includegraphics[width=\linewidth]{figures/fig_retrieval_divergence.png}
  \caption{Per-topic retrieval: cosine similarity (bars) vs.\
  keyword match rate (markers).}
  \label{fig:retrieval}
\end{figure}

\subsection{Ablation Study: Latency vs. Safety with and without the Critic}
\label{sec:ablation}
Same five multi-drug queries, two configurations.
\textbf{Pipeline A (no critic)}: monkey-patch
\texttt{use\_llm\_safety}=\texttt{False}, return the synthesizer
draft as-is. \textbf{Pipeline B}: route every draft through the
dual-agent auditor. Everything else identical.
Wall-clock numbers in Table~\ref{tab:ab}. A: \num{9.50}~s mean,
\num{13.45}~s P95. B: \num{18.55}~s mean, \num{28.80}~s P95. Mean
overhead $\sim$\num{9}~s; P95 overhead \num{15.35}~s, mostly from
critic-triggered regenerations on the heaviest query (Patient
\texttt{10040025}: \num{69} meds, \num{40}+ diagnoses).

Two intercepts from the drug-interaction stress run show what gets
blocked. \emph{Patient} \texttt{10040025}: high-dose ibuprofen plus
amiodarone proposed against a chart already carrying CKD, CHF,
warfarin, digoxin. Pipeline~B refused both, citing those entities
verbatim. \emph{Patient} \texttt{10014354}: warfarin plus
ciprofloxacin proposed against \num{88} active meds. Pipeline~B
surfaced bleeding amplification, QT-prolongation, metformin
renal-dose, and CYP1A2 inhibition, plus
\texttt{[QT-PROLONGATION NOTE]} and \texttt{[RENAL DOSE WARNING]}
locks the single-pass draft missed. Nine seconds per query buys
an audited draft anchored to the chart and verified Hard Lock
surfacing. Clinical context, that's the right side of the
trade-off.


% \subsection{Ablation Study: Latency vs. Safety with and without the Critic}
% \label{sec:ablation}
% To quantify the architectural cost of the Layer~5 adversarial
% critic against the safety value it delivers, we ran the same five
% complex multi-drug queries through two pipeline configurations.
% \textbf{Pipeline A (no critic)} reproduces a standard single-pass
% LLM by monkey-patching the runtime flag
% \texttt{use\_llm\_safety} to \texttt{False}, which returns the
% initial synthesizer draft directly to the user. \textbf{Pipeline B
% (full neuro-symbolic stack)} routes every draft through the
% dual-agent auditor before the response leaves the backend. The
% ablation isolates a single architectural variable; all other
% components (graph retrieval, vector retrieval + cross-encoder
% rerank, Polypharmacy Hard Locks, Visual-RAG) are identical across
% both arms.

% \paragraph*{Latency overhead.} Table~\ref{tab:ab} reports the
% measured wall-clock cost of the critic. Pipeline~A averaged
% \num{9.50}~s mean and \num{13.45}~s P95 across the five queries;
% Pipeline~B averaged \num{18.55}~s mean and \num{28.80}~s P95.
% The mean safety overhead is approximately \num{9}~s; the P95
% overhead reaches \num{15.35}~s and reflects critic-triggered
% regenerations on the most complex query (Patient
% \texttt{10040025}, \num{69} active medications,
% \num{40}+ diagnoses).

% \paragraph*{Safety gain: what the critic intercepted.}
% Two intercepts from the drug-interaction stress run summarize the
% failure mode the critic blocks. \emph{Patient \texttt{10040025}}:
% high-dose ibuprofen plus amiodarone proposed against a chart with
% CKD, CHF, warfarin, and digoxin; Pipeline~B refused both, anchored
% verbatim to those chart entities. \emph{Patient \texttt{10014354}}:
% warfarin plus ciprofloxacin proposed against an \num{88}-medication
% regimen; Pipeline~B surfaced all four interaction axes (bleeding
% amplification, QT-prolongation, metformin renal-dose, CYP1A2
% inhibition) plus the \texttt{[QT-PROLONGATION NOTE]} and
% \texttt{[RENAL DOSE WARNING]} hard locks the single-pass draft
% omitted.

% \paragraph*{Why the trade-off is acceptable.} 
The ablation makes
the trade-off acceptable. Pipeline~A is faster by approximately nine
seconds per query but ships unaudited drafts. Pipeline~B is slower
by the same amount but ships drafts that have been compared
against the patient context by an isolated second LLM and that
have been verified to surface every Hard Lock alert from Layer~4.
In a clinical context, a nine-second delay is preferable to an
unsafe recommendation reaching a physician: the cost of a missed
contraindication on a high-polypharmacy chart materially exceeds
the cost of a longer response time, and the latency budget can be
reduced further in future work
(Section~\ref{sec:future}) without removing the safety property.


\begin{table}[H]
\centering
\caption{A/B latency (5 multi-drug queries). Mean delta from per-query log is \num{9.04}~s; rounded means differ by \num{0.01}~s.}
\label{tab:ab}
\small
\begin{tabular}{@{}lccc@{}}
\toprule
\textbf{Metric} & \textbf{Pipeline A} & \textbf{Pipeline B} & \textbf{Delta}\\
                & \textbf{(no critic)} & \textbf{(critic)} & \\
\midrule
Mean latency & \num{9.50}~s  & \num{18.55}~s & $+\num{9.04}$~s \\
P95 latency  & \num{13.45}~s & \num{28.80}~s & $+\num{15.35}$~s \\
\bottomrule
\end{tabular}
\end{table}

The P95 case is patient \texttt{10040025} (\num{40}+ diagnoses,
\num{69} active meds): Pipeline~B forced a regeneration when the
critic caught an unsafe dosage path. Table~\ref{tab:safety} adds
safety outcomes for the same A/B configurations from the 8-case
adversarial suite.

\begin{table}[H]
\centering
\caption{A/B safety outcomes (8-case adversarial suite). Pipeline~A
counts are counterfactual replays.}
\label{tab:safety}
\small
\begin{tabular}{@{}lcc@{}}
\toprule
\textbf{Outcome (of \num{8})} & \textbf{A (no critic)} & \textbf{B (critic)}\\
\midrule
Unsafe draft reached user & \num{4} & \num{0}\\
Critic audits invoked     & \num{0} & \num{6}\\
Critic overrides          & \num{0} & \num{1}\\
Crash rate                & \num{0}\% & \num{0}\%\\
\bottomrule
\end{tabular}
\end{table}

\begin{figure}[H]
  \centering
  \includegraphics[width=\linewidth]{figures/fig_ablation_latency.png}
  \caption{A/B latency. Pipeline~B stacks the safety overhead on
  Pipeline~A's base.}
  \label{fig:latency}
\end{figure}

\subsection{Adversarial Stress Suite}
\label{sec:adv-results}
Eight stress cases. Two standard consults
(\texttt{10040025}, \texttt{10016742}). One multimodal
contradiction: mandible X-ray \texttt{test\_1800}. One multimodal
pathology with a hidden contraindication: hip MRI
\texttt{test\_1815} plus Patient \texttt{10023239}. One drug-
pharmacokinetics query (\texttt{10002428}). One off-topic prompt
(Star Wars). One direct prompt injection (bleach). One
discharge-summary generation (\texttt{10035631}). Crash rate:
\num{0}~percent. Seven of eight cases scored \num{9}/\num{10}; the
eighth triggered a Layer~8 critic override that blocked the draft
before it reached the user.

\paragraph*{Four-axis breakdown (seven non-override cases).}
\emph{Faithfulness} \num{10}/\num{10} on every case. No clinical
entity in any response was absent from the patient context; that
is the load-bearing claim of the decomposition.
\emph{Completeness} \num{8}--\num{10}. Partial deductions cluster
on dense charts (\texttt{10040025}, \texttt{10002428},
\texttt{10035631}) and reflect the scoped-summarizer design;
chart history irrelevant to the immediate query gets filtered, not
enumerated. \emph{Safety} \num{10}/\num{10} on every case,
including both adversarial prompts. Star Wars returned a refusal
anchored to an empty Layer~2 graph context. Bleach got the same
treatment; the LLM never engaged with the harmful content at all.
\emph{Calibration} \num{7}--\num{9}. Lowest deductions hit cases
where the system declined to extrapolate beyond the imaging
modality. Mandible X-ray scored \num{7} for refusing to search
for pneumonia on a dental panoramic; clinically correct,
judge-conservative. Case~4 (hip MRI + prednisone) sits outside
this scoring because the Layer~8 critic intercepted the draft and
substituted a flagged override.


\subsection{Blind Vision}
\label{sec:vision-results}
Three cases. Table~\ref{tab:vision} has the breakdown. Right
pulmonary artery filling defect: clean localization,
\num{10}/\num{10}. C1 mass: defensive differential rather than a
single cancer commitment, \num{9}/\num{10}. \texttt{test\_1806}:
emphysema and pneumomediastinum matched against the anchor bank,
\num{8}/\num{10}. The \num{8} on \texttt{test\_1806} is full
passes on Faithfulness, Safety, and Pathological Match, minus one
Calibration point for not flagging image-resolution limits.

\begin{table*}[t]
\centering
\caption{Blind Vision suite: per-case radiographic breakdown.}
\label{tab:vision}
\small
\begin{tabular}{@{}lL{0.22\linewidth}L{0.22\linewidth}L{0.22\linewidth}c@{}}
\toprule
\textbf{Test Scan} & \textbf{Structural Localization} & \textbf{Pathological Match} & \textbf{Defensive Hedging} & \textbf{Verdict}\\
\midrule
\texttt{test\_1963} (PE) & Isolated intraluminal defect in right pulmonary artery branch. & Perfectly matched pulmonary embolism against historical anchor bank. & Did not assert beyond visible perfusion defect. & \textbf{10/10}\\
\addlinespace
\texttt{test\_1852} (Neoplasm) & Identified asymmetric soft-tissue swelling at C1 base. & Matched mass-like neoplastic process. & Appropriately deferred definitive cancer diagnosis pending biopsy. & \textbf{9/10}\\
\addlinespace
\texttt{test\_1806} (Emphysema) & Isolated trapped air pockets inside soft tissue structures. & Accurately matched pneumomediastinum/emphysema. & Did not flag image-resolution limits as an explicit hedge. & \textbf{8/10}\\
\bottomrule
\end{tabular}
\end{table*}

\subsection{Golden Dataset Validation}
\label{sec:golden}

\paragraph*{Why precision, not recall, is the safety metric here.}
Recall counts the share of true chart entries reproduced in the
response. The pipeline is a scoped summarizer; a single-drug-pair
query against an \num{88}-medication chart is supposed to return
only the relevant subset. Low recall is the design, not a failure.
Precision (groundedness) is the safety metric: every clinical
entity surfaced has to map to a real chart entry. F1 and safety
recall are reported alongside.

Table~\ref{tab:golden} gives per-patient results across three
identical evaluation runs at temperature~0. Macro precision sits
between \num{89.0} and \num{93.3}~percent across runs (\num{2.0}~pp
spread); F1 lands between \num{23.4} and \num{25.4}~percent.
Patient \texttt{10023239} is the outlier. Run~3 precision drops to
\num{70.0}~percent, traced to a Safety Critic override on Q5 that
reshaped the response format enough for the entity matcher to miss
every medication name. Drop that one query and the spread for the
patient falls under \num{4}~pp, in line with the rest of the suite.
Safety recall averages \num{62}~percent. Precision stays the
load-bearing safety metric because every surfaced entity traces to
a real chart entry; low recall reflects scoped answering.

\paragraph*{What the H column actually counts.} H ranges
\num{18}--\num{26} per run (\num{0.7}--\num{1.0} per query) and
needs a careful read. The Section~\ref{sec:metric-defs}
definition treats a severe medical hallucination as a clinically
actionable claim absent from \texttt{true\_diagnoses} or
\texttt{true\_medications}. H is not that. H counts entities that
the strict Levenshtein matcher flags after the Safety Critic has
already intercepted, annotated, or downgraded the surrounding
claim. Most are class references the chart does not list
verbatim (``NSAID'' on a chart listing \texttt{ibuprofen}; ``beta
blocker'' on a chart listing \texttt{metoprolol}), critic-inserted
hedge tokens (``possible'', ``consider'', ``differential''), and
bold structural artifacts that slip past
\texttt{CLINICAL\_SECTION\_HEADERS} on edge cases. None hit the
severe-hallucination bar. The Safety Critic either reformatted
the claim or appended a flag the matcher then read as a
non-vocabulary token. H therefore reflects the critic's working
surface, not unsafe output reaching clinicians; the zero
severe-hallucination claim in the abstract holds.

\begin{table*}[t]
\centering
\caption{Per-patient Golden Dataset results across three runs
(\num{25} queries, 5 patients; $T{=}0$). Prec./F1 on med+diagnosis
only; vital discrepancies excluded. Safety = mean safety recall.
H = hallucinated entity count (R1/R2/R3).}
\label{tab:golden}
\small
\setlength{\tabcolsep}{7pt}
\begin{tabular*}{\linewidth}{@{\extracolsep{\fill}}
  lc rrr rrr rrr cc @{}}
\toprule
& & \multicolumn{3}{c}{\textbf{Run 1}}
  & \multicolumn{3}{c}{\textbf{Run 2}}
  & \multicolumn{3}{c}{\textbf{Run 3}} & & \\
\cmidrule(lr){3-5}\cmidrule(lr){6-8}\cmidrule(lr){9-11}
\textbf{Patient} & \textbf{Q}
  & \textbf{Prec.} & \textbf{Rec.} & \textbf{F1}
  & \textbf{Prec.} & \textbf{Rec.} & \textbf{F1}
  & \textbf{Prec.} & \textbf{Rec.} & \textbf{F1}
  & \textbf{Safety} & \textbf{H}\\
\midrule
\texttt{10014354} & 5
  & 92.5\% & 11.3\% & 19.6\%
  & 97.5\% & 13.2\% & 22.8\%
  & 97.5\% & 12.4\% & 21.5\%
  & 65\% & 7/3/3\\
\texttt{10002428} & 5
  & 95.7\% & 13.2\% & 22.6\%
  & 93.2\% & 16.0\% & 27.0\%
  & 94.4\% & 16.4\% & 27.0\%
  & 19\% & 3/4/4\\
\texttt{10040025} & 5
  & 92.3\% & 17.4\% & 29.2\%
  & 93.6\% & 14.9\% & 25.2\%
  & 91.2\% & 14.9\% & 25.2\%
  & 100\% & 8/3/4\\
\texttt{10015860} & 5
  & 92.1\% & 13.9\% & 23.7\%
  & 92.1\% & 13.9\% & 23.7\%
  & 92.1\% & 16.5\% & 27.8\%
  & 80\% & 5/5/7\\
\texttt{10023239} & 5
  & 91.1\% & 20.0\% & 31.9\%
  & 90.0\% & 13.8\% & 22.3\%
  & 70.0\% &  9.0\% & 15.4\%\textsuperscript{$\dagger$}
  & 28\% & 3/3/4\\
\midrule
\multicolumn{2}{@{}l}{Macro avg.}
  & 92.7\% & 15.2\% & 25.4\%
  & 93.3\% & 14.3\% & 24.2\%
  & 89.0\% & 13.8\% & 23.4\%
  & 62\% & 26/18/22\\
\multicolumn{2}{@{}l}{Vital discrepancies}
  & \multicolumn{10}{c}{\num{2} (excluded from precision denominator)}\\
\bottomrule
\multicolumn{13}{@{}p{\linewidth}}{%
  \textsuperscript{$\dagger$}Run~3 Q5 (\texttt{10023239}): the
  Safety Critic override changes response formatting; the entity
  matcher finds zero medication matches.  Excluding this query,
  the patient's precision spread is $<$\num{4}~pp.}\\
\end{tabular*}
\end{table*}

%===========================================================
\section{Future Work}
\label{sec:future}
Four directions push the pipeline past its current scope.

\textbf{(1) Imaging beyond pulmonology, plus compound figures.}
The Visual-RAG anchor bank only covers pulmonology today. Too
narrow for production. Next ingestion pass adds cardiology,
radiology, and pathology cohorts. The same pass tackles compound
figures (axial / sagittal / coronal panels laid out in one JPEG),
which the pipeline currently treats as one image and which can
dilute panel-level semantics.

\textbf{(2) More deterministic coverage on drug safety.} A handful
of drug pairs do not belong in the model layer at all. NSAID plus
warfarin, QT-prolonging combinations, CYP1A2 inhibitors plus
narrow-therapeutic-index agents. A deterministic pharmacology
backstop sits upstream of the LLM and removes those classes from
any model decision.

\textbf{(3) Latency: streaming critic and fine-tuning.} The
current critic runs serially after plan generation. A streaming
critic auditing sentence-level chunks in parallel cuts response
time; this needs a backend rewrite. Fine-tuning on clinical
reasoning traces (DeBERTa-v3~\cite{he2021deberta} or a Llama-3.2-1B
distill) should push routine queries close to real-time.

\textbf{(4) Sharper Plan component.} The Plan layer is only as
good as the training signal it gets. Clinical guidelines, FDA
labels, and pharmacist-validated reasoning are the missing inputs;
adding them should improve both accuracy and citation quality.
Production deployment also needs HIPAA-compliant infrastructure
and institutional re-validation
beyond MIMIC-IV Demo and MIMIC-CXR.

%===========================================================
\section{Recommendations}
\label{sec:recommendations}
The following deployment safeguards follow from the evaluation in
Section~\ref{sec:results}. They should be adopted as a single
package; partial adoption weakens the safety guarantees
demonstrated in this paper.

\begin{itemize}
\item \textbf{R1: Keep the safety critic enabled for clinical
queries.} Case~4 of the adversarial suite
(Section~\ref{sec:adv-results}) shows the critic is the only
component that intercepts plausible-but-unsafe drafts; disabling
it reverts the system to the unaudited Pipeline~A baseline.
\item \textbf{R2: Operate as a clinician-in-the-loop assistant.}
The pipeline is a decision-support layer for licensed clinicians,
not an autonomous diagnostic system; no clinician-comparator study
has been run.
\item \textbf{R3: Require PHI redaction or HIPAA-compliant
hosting before any real-clinical deployment.}
\end{itemize}

%===========================================================
\section{Impact}
\label{sec:impact}
Synapse AI augments clinician decision-making rather than replacing
it; every response carries an audit trail naming the matched graph
node, top-reranked passage, figure panel, and any critic
intervention. Off-topic and adversarial prompts are refused by
starving the LLM of grounded context, not by appending safety
phrases. The paper does not claim parity with specialist review;
no clinician-comparator study has been run.

%===========================================================
\section{Conclusion}
An LLM-based clinical assistant can expose its own failure modes
when the architecture pays for the property. Splitting fact
retrieval, drafting, and adversarial review across separate layers
costs about \num{9}~s per query. What it buys, across five suites:
\num{0}~percent crash rate, \num{7}/\num{8} adversarial cases at
\num{9}/\num{10} with one critic override; \num{0.500} mean
semantic similarity and \num{16}/\num{16} conversational rewrite on
the \num{34}-query retrieval benchmark; ${\sim}\num{92}$~percent
macro-precision, macro-F1 ${\sim}\num{24}$~percent, safety recall
\num{62}~percent, and zero severe hallucinations on the
\num{25}-query Golden Dataset; and \num{10}/\num{10},
\num{9}/\num{10}, \num{8}/\num{10} on the three blind-vision
MIMIC-CXR cases. Next milestones: a streaming critic and a
distilled omission classifier, both aimed at compressing the
safety budget without giving up the safety behavior.


\begin{thebibliography}{99}
\small

\bibitem{topol2019deep}
E.~Topol, ``High-performance medicine: the convergence of human and
artificial intelligence,'' \emph{Nature Medicine}, vol.~25, no.~1,
pp.~44--56, 2019.

\bibitem{rajpurkar2022ai}
P.~Rajpurkar, E.~Chen, O.~Banerjee, and E.~Topol, ``AI in health and
medicine,'' \emph{Nature Medicine}, vol.~28, pp.~31--38, 2022.

\bibitem{singhal2023large}
K.~Singhal et al., ``Large language models encode clinical knowledge,''
\emph{Nature}, vol.~620, pp.~172--180, 2023.

\bibitem{moor2023foundation}
M.~Moor et al., ``Foundation models for generalist medical artificial
intelligence,'' \emph{Nature}, vol.~616, pp.~259--265, 2023.

\bibitem{ji2023survey}
Z.~Ji et al., ``Survey of hallucination in natural language
generation,'' \emph{ACM Computing Surveys}, vol.~55, no.~12, 2023.

\bibitem{edge2024local}
D.~Edge et al., ``From local to global: A graph RAG approach to
query-focused summarization,'' arXiv:2404.16130, 2024.

\bibitem{johnson2023mimiciv}
A.~E.~W.~Johnson et al., ``MIMIC-IV, a freely accessible electronic
health record dataset,'' \emph{Scientific Data}, vol.~10, art.~1,
2023.

\bibitem{johnson2019mimiccxr}
A.~E.~W.~Johnson et al., ``MIMIC-CXR: A de-identified publicly
available database of chest radiographs,'' \emph{Scientific Data},
vol.~6, 2019.

\bibitem{pelka2018roco}
O.~Pelka et al., ``Radiology Objects in COntext (ROCO): A multimodal
image dataset,'' in \emph{MICCAI LABELS Workshop}, 2018.

\bibitem{bodenreider2004umls}
O.~Bodenreider, ``The Unified Medical Language System (UMLS):
integrating biomedical terminology,'' \emph{Nucleic Acids Research},
vol.~32, pp.~D267--D270, 2004.

\bibitem{reimers2019sbert}
N.~Reimers and I.~Gurevych, ``Sentence-BERT: Sentence embeddings using
Siamese BERT networks,'' in \emph{EMNLP}, 2019.

\bibitem{lewis2020rag}
P.~Lewis et al., ``Retrieval-augmented generation for knowledge-intensive
NLP tasks,'' in \emph{NeurIPS}, 2020.

\bibitem{medgemma2024}
Google Research, ``MedGemma: open medical models for healthcare
developers,'' Technical Report, 2024.

\bibitem{karabacak2023mirroring}
M.~Karabacak and K.~Margetis, ``Embracing large language models for
medical applications: opportunities and challenges,''
\emph{Cureus}, vol.~15, no.~5, 2023.

\bibitem{he2021deberta}
P.~He, X.~Liu, J.~Gao, and W.~Chen, ``DeBERTa: Decoding-enhanced
BERT with disentangled attention,'' in \emph{ICLR}, 2021.

\bibitem{zhang2024biomedclip}
S.~Zhang et al., ``A multimodal biomedical foundation model trained
from fifteen million image-text pairs (BiomedCLIP),'' arXiv:2303.00915, 2024.

\bibitem{pal2023medhalt}
A.~Pal, L.~K.~Umapathi, and M.~Sankarasubbu, ``Med-HALT:
Medical-domain hallucination test for large language models,'' in
\emph{Conference on Computational Natural Language Learning},
2023.

\bibitem{omiye2023racial}
J.~A.~Omiye, J.~C.~Lester, S.~Spichak, V.~Rotemberg, and R.~Daneshjou,
``Large language models propagate race-based medicine,''
\emph{npj Digital Medicine}, vol.~6, art.~195, 2023.

\bibitem{sheth2023neurosymbolic}
A.~Sheth, K.~Roy, and M.~Gaur, ``Neurosymbolic AI: Why, what, and
how,'' \emph{IEEE Intelligent Systems}, vol.~38, no.~3,
pp.~56--62, 2023.

\bibitem{rotmensch2017graph}
M.~Rotmensch, Y.~Halpern, A.~Tlimat, S.~Horng, and D.~Sontag,
``Learning a health knowledge graph from electronic medical
records,'' \emph{Scientific Reports}, vol.~7, art.~5994, 2017.

\bibitem{zheng2023judge}
L.~Zheng et al., ``Judging LLM-as-a-judge with MT-Bench and Chatbot
Arena,'' in \emph{NeurIPS Datasets and Benchmarks Track}, 2023.

\bibitem{perez2022red}
E.~Perez et al., ``Red teaming language models with language
models,'' in \emph{EMNLP}, 2022.

\end{thebibliography}

%===========================================================
\section*{Code Availability}
Source code, evaluation harness, and golden datasets are open at
\url{https://github.com/Pragyan-dev/Synapse-AI}.

\end{document}
